% Hier könnnen während der Bearbeitung der Arbeit weitere Packages eingebunden werden.

\usepackage{nomencl}
\usepackage{makeidx}

% sorgt dafür das generierter Text auf deutsch ist (Bsp. Inhaltsverzeichnis)
\usepackage[ngerman]{babel}

% für Tabellen
\usepackage{array}

% zusätzliche Mathesymbole
\usepackage{amsmath,amssymb}
\usepackage{bm}

% Calibri esque Font
\usepackage[sfdefault]{carlito}

% FÜr die Aufzählungen
\usepackage{enumitem}

% zum generieren von Literaturverzeichnis
\usepackage{csquotes}
\usepackage{biblatex}
% Robuste Zitier-Wrapper für bewegte Argumente (Captions, etc.)
\DeclareRobustCommand{\citerobust}[1]{\cite{#1}}

% Kapitel Formatierung anpassen
\usepackage{titlesec}
\titleformat{\chapter}[hang]
  {\normalfont\huge\bfseries}
  {\thechapter}               
  {1em}                       
  {}    

% Standard Kopfzeile
\usepackage{fancyhdr}
\fancypagestyle{fancy}{
  \fancyhf{}
  \rhead{\rightmark}
  \cfoot{\thepage}
  \renewcommand{\headrulewidth}{0.4pt}
}

% für keine Header
\fancypagestyle{plain}{
  \fancyhf{}
  \cfoot{\thepage}
  \renewcommand{\headrulewidth}{0pt}
}

% sorgt dafür, dass Leerzeichen hinter parameterlosen Makros nicht als Makroendezeichen interpretiert werden
\usepackage{xspace}

% zum Laden von Grafiken
\usepackage{graphicx}

%Zum färben von Text
\usepackage{xcolor}

% zum setzen von Ankern
\usepackage{hyperref}

% Standardisiertes Abbildungsformat
\usepackage[font=small,labelfont=bf,labelsep=colon]{caption}
\newcommand{\standardfigure}[4][]{%
  \begin{figure}[htbp]
    \centering
    \includegraphics[width=0.85\textwidth]{#2}
    \caption[#3]{#4}
    \label{fig:#1}
  \end{figure}
}